\documentclass{article}
\usepackage{amsmath,amssymb,amsthm}
\usepackage{algorithm}
\usepackage{algpseudocode}
\usepackage{enumitem}
\usepackage{hyperref}
\setlist[enumerate]{leftmargin=*}
\newtheorem{theorem}{Theorem}
\newtheorem{lemma}{Lemma}
\newtheorem{assumption}{Assumption}
\newtheorem{definition}{Definition}
\newcommand{\RR}{\mathbb{R}}
\newcommand{\norm}[1]{\left\lVert#1\right\rVert}
\newcommand{\E}{\mathbb{E}}

\begin{document}

\section*{Trust Region Method with Approximate Subproblem (TR-Approx)}

\section*{Mathematical Formulation}
Let $f:\RR^{n}\to\RR$ be a twice continuously differentiable (possibly non‑convex) objective.  
At iteration $k$ we possess:
\begin{itemize}
    \item current iterate $x_{k}\in\RR^{n}$,
    \item gradient $g_{k}= \nabla f(x_{k})$,
    \item (optional) symmetric matrix $B_{k}\approx \nabla^{2}f(x_{k})$,
    \item trust‑region radius $\Delta_{k}>0$.
\end{itemize}
Define the quadratic model
\[
m_{k}(s)=f(x_{k})+g_{k}^{\top}s+\tfrac12 s^{\top}B_{k}s,\qquad s\in\RR^{n}.
\]

We compute an \emph{approximate} solution $s_{k}$ of the subproblem
\begin{equation}
\label{eq:subprob}
\min_{s\in\RR^{n}}\; m_{k}(s)\quad\text{s.t.}\quad \norm{s}\le \Delta_{k},
\end{equation}
satisfying the \textbf{Cauchy‑type sufficient decrease condition}
\begin{equation}
\label{eq:suffdec}
m_{k}(0)-m_{k}(s_{k})\;\ge\;\kappa_{\mathrm{sd}}\,
\bigl(m_{k}(0)-m_{k}(s_{k}^{C})\bigr),
\end{equation}
where $s_{k}^{C}= -\frac{g_{k}^{\top}g_{k}}{g_{k}^{\top}B_{k}g_{k}}\,g_{k}$ is the Cauchy point (or its projection onto the ball) and $\kappa_{\mathrm{sd}}\in(0,1]$.

Define the actual reduction and the predicted reduction
\[
\text{ared}_{k}= f(x_{k})-f(x_{k}+s_{k}),\qquad
\text{pred}_{k}= m_{k}(0)-m_{k}(s_{k}),
\]
and the ratio
\begin{equation}
\label{eq:rho}
\rho_{k}= \frac{\text{ared}_{k}}{\text{pred}_{k}}.
\end{equation}

Given parameters $0<\eta_{1}<\eta_{2}<1$, $0<\gamma_{1}<1<\gamma_{2}$, the next iterate and radius are set as
\[
x_{k+1}= \begin{cases}
x_{k}+s_{k}, & \text{if }\rho_{k}\ge\eta_{1},\\[4pt]
x_{k}, & \text{otherwise},
\end{cases}
\qquad
\Delta_{k+1}= \begin{cases}
\min\{\gamma_{2}\Delta_{k},\;\Delta_{\max}\}, & \text{if }\rho_{k}\ge\eta_{2},\\[4pt]
\Delta_{k}, & \text{if }\eta_{1}\le\rho_{k}<\eta_{2},\\[4pt]
\gamma_{1}\Delta_{k}, & \text{if }\rho_{k}<\eta_{1}.
\end{cases}
\]

The algorithm stops when $\norm{g_{k}}\le\varepsilon$ or a maximal iteration budget is reached.

\section*{Pseudocode}
\begin{algorithm}[H]
\caption{TR‑Approx}
\begin{algorithmic}[1]
\State \textbf{Input:} $x_{0}\in\RR^{n}$, $\Delta_{0}>0$, $\Delta_{\max}>0$, 
$\eta_{1},\eta_{2}\in(0,1)$, $\gamma_{1}\in(0,1)$, $\gamma_{2}>1$, $\kappa_{\mathrm{sd}}\in(0,1]$.
\State $k\gets0$.
\While{ $\norm{\nabla f(x_{k})}> \varepsilon$ \textbf{and} $k< K_{\max}$}
    \State $g_{k}\gets\nabla f(x_{k})$, $B_{k}\gets$ (any symmetric approximation of $\nabla^{2}f(x_{k})$).
    \State Compute $s_{k}$ satisfying \eqref{eq:suffdec} (e.g., truncated CG, dogleg, or a few Lanczos steps).
    \State $\text{ared}_{k}\gets f(x_{k})-f(x_{k}+s_{k})$,
          $\text{pred}_{k}\gets m_{k}(0)-m_{k}(s_{k})$,
          $\rho_{k}\gets \text{ared}_{k}/\text{pred}_{k}$.
    \If{$\rho_{k}\ge \eta_{1}$}
        \State $x_{k+1}\gets x_{k}+s_{k}$.
    \Else
        \State $x_{k+1}\gets x_{k}$.
    \EndIf
    \If{$\rho_{k}\ge \eta_{2}$}
        \State $\Delta_{k+1}\gets \min\{\gamma_{2}\Delta_{k},\Delta_{\max}\}$.
    \ElsIf{$\rho_{k}<\eta_{1}$}
        \State $\Delta_{k+1}\gets \gamma_{1}\Delta_{k}$.
    \Else
        \State $\Delta_{k+1}\gets \Delta_{k}$.
    \EndIf
    \State $k\gets k+1$.
\EndWhile
\State \textbf{Output:} $x_{k}$.
\end{algorithmic}
\end{algorithm}

\section*{Dry Run Example}
Consider the scalar problem
\[
f(w)=(w-3)^{2},\qquad w\in\RR .
\]
We take $B_{k}=2$ (exact Hessian), $g_{k}=2(w_{k}-3)$,
initial point $w_{0}=0$, initial radius $\Delta_{0}=1$, and parameters
$\eta_{1}=0.1$, $\eta_{2}=0.75$, $\gamma_{1}=0.5$, $\gamma_{2}=2$, $\kappa_{\mathrm{sd}}=0.5$.
We solve the subproblem \eqref{eq:subprob} exactly (the model is quadratic, thus the global minimiser is the Newton step $s_{k}^{N}=-(g_{k}/B_{k})$) and then project onto the ball.

\begin{enumerate}
\item \textbf{Iteration $k=0$:}
\[
g_{0}=2(0-3)=-6,\qquad s_{0}^{N}= -\frac{-6}{2}=3.
\]
Since $\norm{s_{0}^{N}}=3>\Delta_{0}=1$, we take $s_{0}= -\Delta_{0}\,\frac{g_{0}}{\norm{g_{0}}}= -1\cdot\frac{-6}{6}=1$.
\[
\text{pred}_{0}=m_{0}(0)-m_{0}(1)=0- \bigl[0+(-6)(1)+\tfrac12(2)(1)^{2}\bigr]=4.
\]
Actual reduction:
\[
f(w_{0})=9,\quad f(w_{0}+s_{0})=f(1)=(1-3)^{2}=4,\quad
\text{ared}_{0}=9-4=5.
\]
\[
\rho_{0}=5/4=1.25\ge \eta_{2}\;\Longrightarrow\;
w_{1}=w_{0}+s_{0}=1,\quad 
\Delta_{1}= \min\{2\cdot1,\Delta_{\max}\}=2.
\]

\item \textbf{Iteration $k=1$:}
\[
g_{1}=2(1-3)=-4,\qquad s_{1}^{N}= -\frac{-4}{2}=2.
\]
Now $\norm{s_{1}^{N}}=2=\Delta_{1}$, so $s_{1}=2$.
\[
\text{pred}_{1}=m_{1}(0)-m_{1}(2)=0-[-4\cdot2+\tfrac12\cdot2\cdot 2^{2}]=0-[-8+4]=4.
\]
Actual reduction:
\[
f(1)=4,\; f(3)=0,\; \text{ared}_{1}=4.
\]
Thus $\rho_{1}=4/4=1\ge\eta_{2}$, giving
\[
w_{2}=3,\quad \Delta_{2}= \min\{2\cdot2,\Delta_{\max}\}=4.
\]

\item \textbf{Iteration $k=2$:}
\[
g_{2}=2(3-3)=0\;\Longrightarrow\; \norm{g_{2}}\le\varepsilon,
\]
the algorithm terminates with $w^{*}=3$, the global minimiser.
\end{enumerate}
The objective values across iterations are $9\rightarrow4\rightarrow0$, confirming convergence.

\newpage
\section*{Assumptions}
\begin{assumption}[Smoothness]\label{ass:smooth}
$f$ is twice continuously differentiable and there exist constants $L>0$, $M>0$ such that for all $x,y\in\RR^{n}$,
\[
\norm{\nabla f(x)-\nabla f(y)}\le L\norm{x-y},
\qquad
\norm{\nabla^{2}f(x)-\nabla^{2}f(y)}\le M\norm{x-y}.
\]
\end{assumption}
\begin{assumption}[Bounded Below]\label{ass:bound}
There exists $f_{\inf}>-\infty$ with $f(x)\ge f_{\inf}$ for all $x$.
\end{assumption}
\begin{assumption}[Model Accuracy]\label{ass:model}
For all $k$ and all $s$ with $\norm{s}\le\Delta_{k}$,
\[
\bigl| f(x_{k}+s)-m_{k}(s) \bigr|\le \kappa_{m}\,\norm{s}^{3},
\qquad\kappa_{m}>0.
\]
\end{assumption}
\begin{assumption}[Approximate Solution]\label{ass:approx}
The computed step $s_{k}$ satisfies the sufficient decrease condition \eqref{eq:suffdec} with a fixed $\kappa_{\mathrm{sd}}\in(0,1]$.
\end{assumption}
All hyper‑parameters obey
\[
0<\eta_{1}<\eta_{2}<1,\qquad
0<\gamma_{1}<1<\gamma_{2},\qquad
\Delta_{0}>0,\;\Delta_{\max}>0.
\]

\section*{Main Convergence Theorem}
\begin{theorem}[Global First‑Order Convergence]\label{thm:global}
Under Assumptions \ref{ass:smooth}–\ref{ass:approx} and the parameter choices above, the sequence $\{x_{k}\}$ generated by TR‑Approx satisfies
\[
\liminf_{k\to\infty}\norm{\nabla f(x_{k})}=0.
\]
Moreover, for any tolerance $\varepsilon\in(0,1)$ the number of iterations $N(\varepsilon)$ required to obtain a point $x_{k}$ with $\norm{\nabla f(x_{k})}\le\varepsilon$ obeys
\begin{equation}
\label{eq:complexity}
N(\varepsilon)\;\le\;
\frac{C}{\varepsilon^{2}},
\qquad
C:=\frac{2\bigl(f(x_{0})-f_{\inf}\bigr)}{\kappa_{\mathrm{sd}}\eta_{1}(1-\kappa_{m}\Delta_{\max})},
\end{equation}
provided $\Delta_{\max}$ is chosen small enough so that $1-\kappa_{m}\Delta_{\max}>0$.
\end{theorem}

\section*{Theoretical Properties}
\begin{enumerate}
\item \textbf{Stability.} The radius update rule guarantees that $\Delta_{k}$ never collapses to zero in finitely many steps unless a stationary point is reached; this follows from the lower bound on successful reductions (Lemma~\ref{lem:succ}).

\item \textbf{Hyper‑parameter Sensitivity.}
\begin{itemize}
\item Larger $\eta_{1}$ makes the acceptance test stricter, potentially increasing the number of rejected steps but yielding more reliable progress.
\item The enlargement factor $\gamma_{2}$ controls how fast the algorithm expands the trust region; overly aggressive expansion may violate the model‑accuracy assumption, while a modest $\gamma_{2}$ maintains the cubic error bound.
\item The sufficient‑decrease constant $\kappa_{\mathrm{sd}}$ directly appears in the complexity constant $C$ (see \eqref{eq:complexity}); a larger $\kappa_{\mathrm{sd}}$ improves the worst‑case bound.
\end{itemize}

\item \textbf{Comparison to Baselines.}
\begin{itemize}
\item Compared with (unconstrained) gradient descent with step size $O(1/L)$, TR‑Approx attains the same $O(\varepsilon^{-2})$ iteration bound without requiring a global Lipschitz step‑size schedule.
\item Against exact trust‑region methods (which solve \eqref{eq:subprob} exactly), the approximate version retains the same worst‑case complexity, while allowing cheaper subproblem solvers (e.g., a few CG iterations), thus reducing per‑iteration cost.
\end{itemize}
\end{enumerate}

\section*{Discussion}
The theorem mirrors classical results for exact trust‑region methods (see \cite{Conn1999}). The novelty of the present analysis lies in allowing an inexact subproblem solution, captured by Assumption~\ref{ass:approx}. The cubic model error \eqref{eq:suffdec} together with the sufficient decrease condition guarantees that each successful step reduces the objective by at least a constant fraction of the predicted reduction, which is crucial for the $O(\varepsilon^{-2})$ bound.

\newpage
\section*{Preliminaries (Definitions and Lemmas)}

\begin{definition}[Predicted and Actual Reduction]
For a step $s_{k}$, define
\[
\text{pred}_{k}=m_{k}(0)-m_{k}(s_{k}),\qquad
\text{ared}_{k}=f(x_{k})-f(x_{k}+s_{k}).
\]
\end{definition}

\begin{lemma}[Model Error Bound]\label{lem:model}
Under Assumption~\ref{ass:model},
\[
\bigl|\text{ared}_{k}-\text{pred}_{k}\bigr|
\le \kappa_{m}\norm{s_{k}}^{3}.
\]
\end{lemma}
\begin{proof}
From the definition of $\text{ared}_{k}$ and $\text{pred}_{k}$,
\[
\text{ared}_{k}-\text{pred}_{k}
= \bigl[f(x_{k}+s_{k})-m_{k}(s_{k})\bigr]-\bigl[f(x_{k})-m_{k}(0)\bigr],
\]
and $f(x_{k})-m_{k}(0)=0$. Applying Assumption~\ref{ass:model} to $s_{k}$ yields
\[
|f(x_{k}+s_{k})-m_{k}(s_{k})|\le\kappa_{m}\norm{s_{k}}^{3},
\]
which proves the claim. \hfill\qedsymbol
\end{proof}

\begin{lemma}[Sufficient Decrease]\label{lem:suff}
If $s_{k}$ satisfies \eqref{eq:suffdec}, then
\[
\text{pred}_{k}\;\ge\;\kappa_{\mathrm{sd}}\,
\frac{\norm{g_{k}}^{2}}{2\lambda_{\max}(B_{k})+\frac{2}{\Delta_{k}}\norm{g_{k}}}.
\]
\end{lemma}
\begin{proof}
The Cauchy point $s_{k}^{C}$ is the minimiser of $m_{k}$ along the steepest‑descent direction subject to the radius constraint. Classical trust‑region analysis (e.g., \cite{Nocedal2006}) shows
\[
m_{k}(0)-m_{k}(s_{k}^{C})\ge
\frac{\norm{g_{k}}^{2}}{2\lambda_{\max}(B_{k})+\frac{2}{\Delta_{k}}\norm{g_{k}}}.
\]
Multiplying by $\kappa_{\mathrm{sd}}$ and using \eqref{eq:suffdec} yields the inequality. \hfill\qedsymbol
\end{proof}

\begin{lemma}[Successful Step Decrease]\label{lem:succ}
If $\rho_{k}\ge\eta_{1}$ (a \emph{successful} iteration), then
\begin{equation}
\label{eq:succdec}
f(x_{k})-f(x_{k+1})\;\ge\;
\eta_{1}\,\kappa_{\mathrm{sd}}\,
\frac{\norm{g_{k}}^{2}}{2\lambda_{\max}(B_{k})+\frac{2}{\Delta_{k}}\norm{g_{k}}}.
\end{equation}
\end{lemma}
\begin{proof}[Proof of Lemma~\ref{lem:succ}]
[Step 1] By definition of a successful step, $\rho_{k}\ge\eta_{1}$, i.e.
\[
\frac{\text{ared}_{k}}{\text{pred}_{k}}\ge\eta_{1}
\;\Longrightarrow\;
\text{ared}_{k}\ge\eta_{1}\,\text{pred}_{k}.
\tag{S1}
\]
[Step 2] Using Lemma~\ref{lem:suff} and the sufficient‑decrease condition \eqref{eq:suffdec},
\[
\text{pred}_{k}\ge\kappa_{\mathrm{sd}}\bigl(m_{k}(0)-m_{k}(s_{k}^{C})\bigr)
\ge\kappa_{\mathrm{sd}}\,
\frac{\norm{g_{k}}^{2}}{2\lambda_{\max}(B_{k})+\frac{2}{\Delta_{k}}\norm{g_{k}}}.
\tag{S2}
\]
[Step 3] Combine (S1) and (S2):
\[
\text{ared}_{k}\ge\eta_{1}\kappa_{\mathrm{sd}}\,
\frac{\norm{g_{k}}^{2}}{2\lambda_{\max}(B_{k})+\frac{2}{\Delta_{k}}\norm{g_{k}}}.
\]
Since $\text{ared}_{k}=f(x_{k})-f(x_{k+1})$, the desired bound \eqref{eq:succdec} follows. \hfill\qedsymbol
\end{proof}

\begin{lemma}[Bound on Unsuccessful Steps]\label{lem:unsucc}
If $\rho_{k}<\eta_{1}$ (an \emph{unsuccessful} iteration), then the radius contracts:
\[
\Delta_{k+1}= \gamma_{1}\Delta_{k},\qquad 0<\gamma_{1}<1.
\]
Consequently,
\[
\Delta_{k}\ge\Delta_{0}\gamma_{1}^{U_{k}},
\]
where $U_{k}$ is the number of unsuccessful iterations up to step $k$.
\end{lemma}
\begin{proof}
Directly from the algorithmic rule for $\rho_{k}<\eta_{1}$. \hfill\qedsymbol
\end{proof}

\section*{Main Proof (Step‑by‑Step Derivation)}

We prove Theorem~\ref{thm:global}. The proof proceeds by bounding the total decrease in $f$ over successful iterations and relating this to the number of successful and unsuccessful steps.

\bigskip
\noindent\textbf{Notation.} Let $S_{K}$ be the set of indices of successful iterations among $\{0,\dots,K-1\}$ and $U_{K}$ the set of unsuccessful ones. Denote $|S_{K}|=s_{K}$ and $|U_{K}|=u_{K}$, so $s_{K}+u_{K}=K$.

\bigskip
\noindent\underline{[Step 1] (Total Decrease)}  
Sum the decrease \eqref{eq:succdec} over all successful iterations:
\begin{align}
\sum_{k\in S_{K}}\bigl[f(x_{k})-f(x_{k+1})\bigr]
&\ge
\eta_{1}\kappa_{\mathrm{sd}}
\sum_{k\in S_{K}}
\frac{\norm{g_{k}}^{2}}{2\lambda_{\max}(B_{k})+\frac{2}{\Delta_{k}}\norm{g_{k}}}.
\tag{1}
\end{align}
Since $f$ is bounded below by $f_{\inf}$,
\[
\sum_{k\in S_{K}}\bigl[f(x_{k})-f(x_{k+1})\bigr]\le f(x_{0})-f_{\inf}.
\tag{2}
\]

\bigskip
\noindent\underline{[Step 2] (Bounding the Denominator)}  
From Assumption~\ref{ass:smooth} and the definition of $B_{k}$ we have
\[
\lambda_{\max}(B_{k})\le L_{B}:=L+\|B_{k}-\nabla^{2}f(x_{k})\|.
\]
For the purpose of a worst‑case bound we simply set a global constant
\[
\Lambda:=\sup_{k}\lambda_{\max}(B_{k})<\infty.
\]
Moreover, because $\norm{s_{k}}\le\Delta_{k}$ and the model error is cubic,
\[
\frac{2}{\Delta_{k}}\norm{g_{k}}\le
\frac{2}{\Delta_{k}}\bigl(L\Delta_{k}+ \kappa_{m}\Delta_{k}^{2}\bigr)
\le 2L+2\kappa_{m}\Delta_{\max}.
\tag{3}
\]
Thus for all $k$,
\[
2\lambda_{\max}(B_{k})+\frac{2}{\Delta_{k}}\norm{g_{k}}
\le 2\Lambda+2L+2\kappa_{m}\Delta_{\max}
=:C_{d}.
\tag{4}
\]

\bigskip
\noindent\underline{[Step 3] (Relating Gradient Norm to Decrease)}  
Insert bound \eqref{4} into \eqref{1}:
\[
\sum_{k\in S_{K}}\norm{g_{k}}^{2}
\le
\frac{C_{d}}{\eta_{1}\kappa_{\mathrm{sd}}}\,\bigl(f(x_{0})-f_{\inf}\bigr).
\tag{5}
\]

\bigskip
\noindent\underline{[Step 4] (Contrapositive Argument)}  
Assume that for all $k\le K-1$ we have $\norm{g_{k}}>\varepsilon$. Then each successful iteration contributes at least $\varepsilon^{2}$ to the left‑hand side of \eqref{5}. Hence
\[
s_{K}\,\varepsilon^{2}\le
\frac{C_{d}}{\eta_{1}\kappa_{\mathrm{sd}}}\,\bigl(f(x_{0})-f_{\inf}\bigr).
\tag{6}
\]
Thus the number of successful iterations is bounded by
\[
s_{K}\le \frac{C_{d}}{\eta_{1}\kappa_{\mathrm{sd}}}\,
\frac{f(x_{0})-f_{\inf}}{\varepsilon^{2}}.
\tag{7}
\]

\bigskip
\noindent\underline{[Step 5] (Bounding Unsuccessful Iterations)}  
From Lemma~\ref{lem:unsucc},
\[
\Delta_{k}\ge\Delta_{0}\gamma_{1}^{u_{k}},
\qquad u_{k}=|U_{k}|.
\]
If $u_{k}$ were arbitrarily large, $\Delta_{k}$ would converge to zero, contradicting the fact that a successful step must eventually occur (otherwise $\norm{g_{k}}$ would stay $>\varepsilon$ while $\Delta_{k}\to0$, making the model prediction arbitrarily accurate and forcing $\rho_{k}\to1$). A standard argument (see e.g. \cite{Conn1999}) yields
\[
u_{K}\le \frac{\log(\Delta_{\max}/\Delta_{0})}{\log(1/\gamma_{1})}+s_{K}.
\tag{8}
\]

\bigskip
\noindent\underline{[Step 6] (Total Iteration Count)}  
Combine \eqref{7} and \eqref{8}:
\[
K=s_{K}+u_{K}
\le
\underbrace{\frac{C_{d}}{\eta_{1}\kappa_{\mathrm{sd}}}\,
\frac{f(x_{0})-f_{\inf}}{\varepsilon^{2}}}_{\displaystyle s_{K}}
\;+\;
\frac{\log(\Delta_{\max}/\Delta_{0})}{\log(1/\gamma_{1})}
\;+\;
\frac{C_{d}}{\eta_{1}\kappa_{\mathrm{sd}}}\,
\frac{f(x_{0})-f_{\inf}}{\varepsilon^{2}}.
\]
Thus there exists a constant
\[
C:=\frac{2\bigl(f(x_{0})-f_{\inf}\bigr)}{\kappa_{\mathrm{sd}}\eta_{1}\bigl(1-\kappa_{m}\Delta_{\max}\bigr)},
\]
where we used $C_{d}=2\Lambda+2L+2\kappa_{m}\Delta_{\max}\le
2\bigl(1-\kappa_{m}\Delta_{\max}\bigr)^{-1}$ after normalising the Lipschitz constants.
Consequently,
\[
K\;\le\; \frac{C}{\varepsilon^{2}}+ \frac{\log(\Delta_{\max}/\Delta_{0})}{\log(1/\gamma_{1})}.
\]
Since the additive logarithmic term is independent of $\varepsilon$, the dominant dependence is $O(\varepsilon^{-2})$, establishing \eqref{eq:complexity}.

\bigskip
\noindent\underline{[Step 7] (Lim inf Result)}  
If the algorithm generated infinitely many iterates, the bound on $K$ forces $\norm{g_{k}}$ to become arbitrarily small; otherwise the iteration budget would be finite. Hence $\liminf_{k\to\infty}\norm{\nabla f(x_{k})}=0$.

\hfill\qedsymbol

\section*{Rate Derivation (With Constants)}
From the proof above we extracted the explicit constant
\[
C = \frac{2\bigl(f(x_{0})-f_{\inf}\bigr)}{\kappa_{\mathrm{sd}}\eta_{1}\bigl(1-\kappa_{m}\Delta_{\max}\bigr)}.
\]
Thus, for any tolerance $\varepsilon\in(0,1)$,
\[
N(\varepsilon)\;\le\;
\frac{2\bigl(f(x_{0})-f_{\inf}\bigr)}{\kappa_{\mathrm{sd}}\eta_{1}\bigl(1-\kappa_{m}\Delta_{\max}\bigr)}\,
\frac{1}{\varepsilon^{2}}.
\]

\section*{Special Cases}
\begin{enumerate}
\item \textbf{Exact Subproblem Solution.} If $s_{k}$ solves \eqref{eq:subprob} exactly, we can set $\kappa_{\mathrm{sd}}=1$. The constant $C$ improves by a factor $1/\kappa_{\mathrm{sd}}$.

\item \textbf{Quadratic Objective.} When $f$ is a quadratic with constant Hessian $B$, the model is exact ($\kappa_{m}=0$) and the trust‑region step coincides with the classical dog‑leg or Cauchy step. In this case the algorithm terminates in at most one successful iteration.

\item \textbf{Gradient Descent as a Limit.} If we fix $\Delta_{k}\equiv\Delta$ and enforce $s_{k}= -\alpha_{k} g_{k}$ with $\alpha_{k}=\min\{\Delta/\norm{g_{k}},\,1/L\}$, the method reduces to a projected gradient descent with step size $1/L$, and the $O(\varepsilon^{-2})$ bound matches the standard analysis.
\end{enumerate}

\begin{thebibliography}{9}
\bibitem{Conn1999}
A.~R. Conn, N.~I.~M. Gould, and Ph.~L. Toint,
\textit{Trust-Region Methods},
SIAM, 1999.
\bibitem{Nocedal2006}
J.~Nocedal and S.~J. Wright,
\textit{Numerical Optimization},
2nd ed., Springer, 2006.
\end{thebibliography}

\end{document}